\documentclass[11pt,a4paper]{article}
\usepackage[utf8]{inputenc}
\usepackage{geometry}
\usepackage{booktabs}
\usepackage{listings}
\usepackage{xcolor}
\usepackage{hyperref}

\geometry{margin=1in}

\definecolor{codegreen}{rgb}{0,0.6,0}
\definecolor{codegray}{rgb}{0.5,0.5,0.5}
\definecolor{codepurple}{rgb}{0.58,0,0.82}
\definecolor{backcolour}{rgb}{0.95,0.95,0.92}

\lstdefinestyle{mystyle}{
    backgroundcolor=\color{backcolour},   
    commentstyle=\color{codegreen},
    keywordstyle=\color{purple},
    numberstyle=\tiny\color{codegray},
    stringstyle=\color{codepurple},
    basicstyle=\ttfamily\footnotesize,
    breakatwhitespace=false,         
    breaklines=true,                 
    captionpos=b,                    
    keepspaces=true,                 
    numbers=left,                    
    numbersep=5pt,                  
    showspaces=false,                
    showstringspaces=false,
    showtabs=false,                  
    tabsize=2
}
\lstset{style=mystyle}

\title{Cryptography \& Network Security Integrated Situation Report\\ 
\large Module: ETTCS801}
\author{Noela Minicha Angelo Geri \\ \large Roll Number: 4202670037}
\date{September 28, 2026}

\begin{document}

\maketitle

\section{Introduction}
This technical report details the risks, controls, and cryptographic solutions implemented for the ULK Polytechnic Institute central server infrastructure. It outlines our risk management process, python security tool development, and network traffic filtering tests.

\section{Risk Assessment and Rankings}
A security assessment revealed multiple weaknesses in the server architecture. Below are the identified issues, evaluation metrics, and remediation controls.

\subsection{Identified Assets and Vulnerabilities}
\begin{itemize}
    \item \textbf{Data in Transit:} Student data moving between campuses is unencrypted, presenting a high risk of eavesdropping.
    \item \textbf{Central Server Resources:} Exposed management access paired with weak staff passwords leaves the platform open to dictionary attacks.
    \item \textbf{Network Perimeter Control:} Open network paths allow unauthenticated guest Wi-Fi clients direct communication access to the secure server database.
\end{itemize}

\subsection{Risk Rating Matrix}
\begin{table}[h!]
\centering
\begin{tabular}{lllll}
\toprule
Risk Threat Source & Likelihood & Impact & Severity Priority & Reason Justification \\
\midrule
Weak Passwords / Probing & High & High & Critical & Active login attempts are occurring. \\
Guest Network Exposure & High & Medium & High & Open internal routes exist today. \\
Unencrypted Transfers & Medium & High & High & Leads to systemic identity leakage. \\
\bottomrule
\end{tabular}
\caption{ULK Server Infrastructure Risk Rankings}
\label{tab:risks}
\end{table}

\subsection{Recommended Controls}
To mitigate these exposures, we recommend implementing: Transport Layer Security (TLS/HTTPS) for data in transit; Multi-Factor Authentication (MFA) with complex password requirements for staff accounts; and strict Network Segmentation using firewall policies to block guest traffic.

\section{Cryptography and File Integrity Application}
To secure data at rest and verify authenticity, a Python script utility named \texttt{toolkit.py} was developed using the Python \texttt{cryptography} ecosystem library.

\subsection{Functional Implementation}
The software performs symmetric encryption through the \texttt{Fernet} standard. Keys are kept separate from the code repository using a local storage file (\texttt{secret.key}). Integrity checking is handled by hashing the dataset using the \textbf{SHA-256 algorithm}.

\subsection{Execution Verification Outputs}
When run locally, the program yields the following output log verifying success:
\begin{lstlisting}[language=bash]
--- ULK Security Toolkit ---
student_records.csv not found. Creating a sample file...
Original File SHA-256: 8a48efc63c4e...
Successfully encrypted student_records.csv into student_records.enc
Successfully decrypted student_records.enc into student_records_decrypted.csv
Decrypted File SHA-256: 8a48efc63c4e...
Success: The original file and decrypted file match exactly! Integrity verified.
\end{lstlisting}

\section{Network Traffic Filtering Configuration}
A perimeter traffic architecture policy was established to limit communication routes to the server node based on internal network interfaces.

\subsection{Applied Firewall Rules}
The dynamic interface rules below were appended to the system deployment layout:
\begin{lstlisting}[language=bash]
# Block all guest network interface access to the server environment
sudo iptables -A INPUT -i $GUEST_IFACE -j DROP

# Permit authorized staff access to the secure service port (443)
sudo iptables -A INPUT -i $STAFF_IFACE -p tcp --dport 443 -j ACCEPT

# Block remaining inbound access to port 443 from any other interface
sudo iptables -A INPUT -p tcp --dport 443 -j DROP
\end{lstlisting}

\subsection{Verification Test Logs}
Three connection profiles were evaluated to ensure correctness:
\begin{enumerate}
    \item \textbf{Staff to Port 443 (Permitted):} Running \texttt{curl -I https://records-server.local} yielded an active response of \texttt{HTTP/1.1 200 OK}.
    \item \textbf{Guest Network Connection (Blocked):} Running \texttt{ping -c 3 records-server.local} resulted in immediate drops with \texttt{100\% packet loss}.
    \item \textbf{External Connection to Port 443 (Blocked):} Running an untrusted network check timed out successfully, displaying \texttt{curl: (28) Connection timed out}.
\end{enumerate}

\section{Conclusion and Submission Information}
The vulnerabilities highlighted by the initial server security review—specifically weak staff authentication, unsegmented guest paths, and unencrypted cross-campus communications—have been systematically mitigated. By combining robust symmetric encryption (Fernet standard) and SHA-256 integrity verification with strategic interface-based Netfilter/iptables firewall boundaries, the central student records platform has been hardened against both internal guest network exposures and active external probing. 

The complete version-controlled implementation project repository, including code, test logs, and configuration files, can be verified online at the following location:

\begin{center}
\url{https://github.com/NoelaMiniG/cryptography-network-security-exam}
\end{center}

\end{document}
